\section{Wstęp}
\label{sec:wstep}

\subsection{Wprowadzenie}

Trudno dziś wyobrazić sobie firmę, która nie korzysta z chmury obliczeniowej.
Poczta, dokumenty, bazy danych, całe aplikacje --- wszystko to coraz częściej
działa na serwerach wynajmowanych od dostawców takich jak Amazon, Microsoft czy
Google. Razem z usługami do chmury trafiają jednak dane, w tym te najbardziej
wrażliwe. To czyni chmurę łakomym kąskiem dla cyberprzestępców, a \emph{naruszenie
danych} --- czyli ich wyciek lub kradzież --- jednym z najbardziej kosztownych
incydentów, jakie mogą spotkać organizację.

Klasyczne zabezpieczenia, oparte na gotowych ,,sygnaturach'' znanych ataków,
coraz słabiej radzą sobie z nowymi zagrożeniami i z ogromną skalą ruchu w
chmurze. Dlatego dużą nadzieję pokłada się w \emph{uczeniu maszynowym} --- czyli
w modelach, które samodzielnie uczą się rozpoznawać złośliwą aktywność na
podstawie wcześniejszych przykładów. Wśród wielu dostępnych metod szczególnie
wyróżniają się dwie: \emph{lasy losowe} (sprawdzona, ,,klasyczna'' technika) oraz
\emph{uczenie głębokie} (nowsze, oparte na sieciach neuronowych). Obie potrafią
osiągać świetne wyniki, ale różnią się złożonością, kosztem i tym, jak łatwo
zrozumieć ich decyzje.

\subsection{Cel pracy}

Celem pracy jest \textbf{porównanie skuteczności metod uczenia głębokiego i
lasów losowych w wykrywaniu (prognozowaniu) naruszeń danych w chmurze}. Przez
,,prognozowanie naruszenia'' rozumiemy rozpoznawanie pojedynczych połączeń
sieciowych jako normalnych lub złośliwych --- wykrycie podejrzanego połączenia
jest bowiem wczesnym ostrzeżeniem przed pełnym wyciekiem danych.

Chcemy w szczególności odpowiedzieć na trzy pytania:
\begin{enumerate}
  \item Które podejście --- las losowy czy sieci głębokie --- skuteczniej
        wykrywa ataki?
  \item Ile kosztuje (pod względem czasu uczenia i działania) każde z nich i czy
        wyższa skuteczność jest warta tego kosztu?
  \item Czy bardziej rozbudowane sieci (CNN, LSTM) są lepsze od prostej sieci
        MLP na danych opisujących ruch sieciowy?
\end{enumerate}

\subsection{Zakres i sposób realizacji}

Projekt ma charakter zarówno teoretyczny, jak i praktyczny. Część praktyczna to
napisany od podstaw program w Pythonie, który realizuje cały eksperyment ---
od przygotowania danych, przez uczenie pięciu modeli, po obliczenie wyników i
wygenerowanie wykresów. Wszystkie liczby i rysunki w sprawozdaniu pochodzą
bezpośrednio z tego programu, dzięki czemu treść jest zgodna z faktycznymi
wynikami.

\subsection{Układ dokumentu}

Dalsza część pracy ma typowy układ: w rozdziale~\ref{sec:problem} przybliżamy
problem bezpieczeństwa danych w chmurze, w rozdziale~\ref{sec:podstawy}
opisujemy porównywane metody, rozdział~\ref{sec:metodyka} przedstawia sposób
przeprowadzenia eksperymentu, rozdział~\ref{sec:realizacja} prezentuje i omawia
wyniki, a rozdział~\ref{sec:wnioski} zawiera wnioski oraz pomysły na dalsze
prace.
